% !TeX root = ../main.tex
% Add the above to each chapter to make compiling the PDF easier in some editors.

\chapter{Related Work}\label{chapter:related-work}

\section{Iterative Solvers and Preconditioners}\label{sec:solvers-background}

Iterative methods for sparse linear systems construct a sequence of approximations to the
solution of $Ax = b$ by building a Krylov subspace
$\mathcal{K}_k(A, r_0) = \operatorname{span}\{r_0,\, Ar_0,\, \ldots,\, A^{k-1}r_0\}$
from the initial residual $r_0$.
The Conjugate Gradient method (CG) is the standard choice for symmetric positive definite
systems, while methods such as GMRES, FGMRES, DGMRES, BiCGStab, and BiCGStabl handle
non-symmetric or indefinite problems.
Further specialisations include MINRES for symmetric indefinite systems, as well as CR and
CGS for specific structural properties.
A comprehensive treatment of Krylov subspace methods is given by Saad~\cite{saad2003iterative}.

Preconditioners transform the linear system to improve convergence by approximating
$A^{-1}$ at low computational cost.
Simple diagonal methods such as Jacobi, SOR, and the Eisenstat variant are inexpensive but
effective for diagonally dominant systems.
Incomplete LU factorisation (ILU), introduced by Meijerink and Van der
Vorst~\cite{meijerink1977iterative}, approximates the LU decomposition while discarding
small fill-in entries, offering a stronger approximation at moderate cost.
Block variants such as block Jacobi and additive Schwarz (ASM) decompose the system into
local subproblems that can be solved independently or in parallel.
Algebraic multigrid (AMG and its PETSc implementation GAMG) constructs a hierarchy of
coarser representations and achieves near-optimal convergence for a wide class of
problems~\cite{stuben2001review}.
Because no single solver-preconditioner pair dominates across all matrix types, selecting
the right combination requires problem-specific knowledge, which motivates automated
selection.

\section{Automated Solver Selection}\label{sec:automated-selection}

Automating the choice of solver and preconditioner has been studied from several angles.
Rule-based expert systems encode practitioner knowledge about which configurations suit
which matrix properties, but are limited to known problem classes and require significant
manual effort to maintain.
Classical machine learning approaches overcome this by training classifiers such as support
vector machines or random forests on hand-crafted matrix statistics, learning the mapping
from matrix properties to optimal solver configurations directly from benchmark data.

More recent work extends this idea to richer input representations.
MM-AutoSolver~\cite{xiong2025mmautosolver} supplements scalar matrix features with a CNN
branch processing a visual rendering of the sparsity pattern, improving accuracy over
feature-only baselines.
Weng et al.~\cite{weng2026embedding} take a different direction with an embedding-based
framework that learns solver-problem relationships directly from observed performance data,
achieving substantial improvements over classical feature-based models across a broad set
of PETSc solver-preconditioner configurations.

\section{MM-AutoSolver}\label{sec:mm-autosolver}

MM-AutoSolver~\cite{xiong2025mmautosolver} is the primary baseline and architectural
reference for this thesis.
The model combines two input streams: a scalar feature branch processing 17 hand-crafted
matrix statistics through an MLP, and a CNN branch operating on a $128 \times 128$ pixel
density image of the sparse matrix.
The two branches are fused via element-wise addition before a shared classification head
predicts one of 19 fixed solver-preconditioner pairs from PETSc~\cite{dalcin2011parallel}.
The CNN uses two convolutional layers with Tanh activations and $3 \times 3$ and
$5 \times 5$ kernels respectively, trained with Adam ($\text{lr} = 10^{-3}$,
256 epochs, batch size 512).
On a dataset of 11{,}623 matrices from the SuiteSparse collection~\cite{davis2011university} the model achieves Acc\,=\,78.54\,\%,
MP\,=\,63.41\,\%, MR\,=\,62.81\,\%, and F1\,=\,62.53\,\%.
This thesis independently replicates MM-AutoSolver on a custom dataset~(\autoref{sec:mm-replication})
and extends the approach by systematically evaluating 14 different image encoding strategies,
dual-channel CNN inputs, and ensemble methods.

\section{Embedding-based Solver Performance Prediction}\label{sec:supervisor-framework}

Weng et al.~\cite{weng2026embedding} present an embedding-based framework for linear solver
performance prediction evaluated across 101 PETSc solver-preconditioner configurations on
a broad set of SuiteSparse benchmark matrices.
Rather than engineering features by hand, the framework learns solver-problem relationships
directly from observed performance data and projects unseen problems into a learned
embedding space using inexpensive numerical features, achieving a 37\,\% reduction in mean
average percentage error and a 17\,\% improvement in top-prediction accuracy over classical
feature-based models.
This thesis targets a subset of 19 of the 101 solver-preconditioner configurations from
this framework as classification labels, with the goal of training a CNN-based model that
could eventually serve as a lightweight inference module within the broader
evaluation pipeline.

\section{Image Representations of Sparse Matrices}\label{sec:image-representations}

Representing a sparse matrix as a two-dimensional image allows a CNN to exploit spatial
and structural patterns that scalar statistics cannot capture.
The simplest encoding is a binary sparsity pattern, marking each non-zero entry as a pixel.
Denser representations encode the value magnitude or local non-zero density per pixel cell,
optionally on a logarithmic scale to compress dynamic range; signed variants additionally
preserve entry signs.
Before rendering, the Reverse Cuthill--McKee (RCM) reordering~\cite{cuthill1969reducing}
reduces the matrix bandwidth by permuting rows and columns, concentrating non-zero entries
closer to the diagonal and making structural patterns more visually compact.
MM-AutoSolver~\cite{xiong2025mmautosolver} uses a single density encoding at
$128 \times 128$ pixels; this thesis evaluates 14 distinct encoding strategies at two
resolutions to identify the most informative representation for solver selection.

\section{Convolutional Neural Networks and Multimodal Learning}\label{sec:cnn-background}

Convolutional neural networks learn hierarchical spatial features from image data and are
well established for image classification tasks.
In this thesis, the CNN serves as one branch of a multimodal architecture that combines
image and scalar inputs.
Multimodal machine learning, as surveyed by Baltru\v{s}aitis et
al.~\cite{baltruvsaitis2018multimodal}, combines complementary signals from multiple input
modalities to improve prediction over any single modality alone.
The dual-stream design used here follows this paradigm: the CNN branch captures spatial
matrix structure while the MLP branch processes scalar statistics, and the two sources of
information are fused through concatenation before a shared classification head.
